\chapter{Analytická časť}\label{ch:analytics}

Analytická časť práce vytvára teoretický a technický základ pre návrh a implementáciu riešenia. Jej cieľom je systematicky analyzovať existujúci stav, identifikovať jeho obmedzenia a na základe toho navrhnúť vhodnú architektúru umožňujúcu porovnanie monolitického a mikroservisného prístupu.

\subsection{Úvod do analytickej časti}

Analytická časť predstavuje detailnú technickú analýzu problému, aktuálneho riešenia a možností jeho optimalizácie prostredníctvom alternatívnej architektúry. Cieľom je identifikovať limity existujúcej monolitickej implementácie, definovať funkčné a nefunkčné požiadavky, preskúmať vhodné technológie a navrhnúť mikroservisnú architektúru, ktorá umožní objektívne porovnanie oboch prístupov.

\subsection{Analýza technologického kontextu}

Moderné webové aplikácie využívajú škálovateľné architektúry, automatizované nasadzovanie a kontajnerové riešenia. V tejto práci sa preto používa technologický stack založený na frameworku \gls{springboot} pre tvorbu \gls{restapi}, relačnej databáze \gls{postgresql} pre perzistenciu dát, knižnici \gls{react} pre implementáciu single-page aplikácie a nástrojoch \gls{docker} a \gls{dockercompose} pre kontajnerizáciu a reprodukovateľné vývojové prostredie. APIs sú centralizované cez \gls{springcloudgateway}, ktorý poskytuje jednotný vstupný bod a základné bezpečnostné a smerovacie mechanizmy. Tento stack je vhodný pre obe architektúry – \gls{monolith} aj \gls{microservice} – a umožňuje ich férové porovnanie v rovnakých technologických podmienkach.

\subsection{Analýza existujúcej monolitickej implementácie}

V tejto podkapitole je analyzovaná existujúca monolitická implementácia aplikácie hry Dáma. Analýza sa zameriava na architektúru systému, spôsob správy stavu, internú komunikáciu medzi komponentmi a prevádzkové vlastnosti riešenia. Cieľom je identifikovať silné stránky monolitu a zároveň jeho limity z pohľadu škálovateľnosti a dlhodobej udržateľnosti.

\subsubsection{Architektúra}

Monolitická verzia pozostáva z jedného backendového modulu implementovaného v Spring Boot~3. Aplikácia obsahuje vrstvu kontrolérov, vrstvu služieb pre spracovanie hernej logiky (validácia ťahov, povinné bije, kráľovanie), vrstvu pre ukladanie dát pomocou Spring Data JPA a databázy PostgreSQL. Frontend je samostatná React SPA aplikácia komunikujúca cez REST API.

\subsubsection{Stav a interná komunikácia}

Všetky komponenty zdieľajú ten istý runtime a databázu. To umožňuje veľmi rýchlu internú komunikáciu, jednoduché testovanie a minimálne prevádzkové nároky, no zároveň vytvára silné prepojenie medzi modulmi.

\subsubsection{Popis existujúceho riešenia}

Existujúce riešenie pozostáva z monolitickej aplikácie, ku ktorej je v rámci práce navrhnutá a implementovaná aj mikroservisná verzia. Obe verzie používajú frontend implementovaný v Reacte a backend postavený na Spring Boote s databázou PostgreSQL. Všetky údaje (nicky, avatary, priebeh hry, výsledky, komentáre a hodnotenia) sa načítavajú cez HTTP požiadavky na backend. Interakcia prebieha medzi dvoma hráčmi na jednom zariadení, pričom frontend aplikácia v Reacte zabezpečuje celú vizuálnu časť.

\subsubsection{Prevádzkové správanie a vlastnosti}

\begin{itemize}
\item práca s používateľskými údajmi: frontend načítava všetky potrebné údaje (nicky, avatary, výsledky) cez REST API;
\item perzistencia výsledkov: výsledky sú ukladané cez služby monolitu do PostgreSQL;
\item jednovláknová konzistencia: všetky operácie bežia v jednom backendovom procese;
\item Dev/Ops zjednodušenie: monolit je ľahko spustiteľný lokálne – samostatný backend a samostatný frontend.
\end{itemize}

\subsubsection{Identifikované limity a problémy}

\begin{itemize}
\item škálovanie: pri náraste zaťaženia (viac súbežných relácií) je potrebné škálovať celý monolit, čo je neefektívne pri potrebe zvýšiť kapacitu len časti funkcionality;
\item odolnosť voči chybám: chyba v jednom module môže spôsobiť zlyhanie celej aplikácie;
\item vývojová súbežnosť: viacerí vývojári pracujúci na rôznych častiach môžu naraziť na konflikty v kóde a nasadzovaní;
\item testovateľnosť: je ťažšie izolovať moduly pre samostatné testovanie, integrácia je náročnejšia;
\item nasadzovanie: každá zmena vyžaduje redeploy celého systému;
\item dlhodobá údržba komplikuje pridávanie nových funkcií, keďže zmeny v jednej časti môžu mať široký dopad;
\item monolitická štruktúra obmedzuje možnosť samostatného rozvoja vybraných domén (napr. len správa hráčov alebo skóre).
\end{itemize}

Napriek uvedeným obmedzeniam zostáva monolitická implementácia dôležitým referenčným bodom pre celé porovnanie. Predstavuje prirodzený „východiskový stav“, ktorý by vznikol aj v mnohých reálnych projektoch, kde je hlavným cieľom čo najrýchlejšie dodať funkčný prototyp. V rámci tejto práce je monolitické riešenie implementované tak, aby bolo čo najčistejšie z pohľadu architektonického členenia (vrstvy kontrolérov, služieb a perzistencie), no zároveň zostáva typickým predstaviteľom jednoduchšej architektúry s jedným nasaditeľným artefaktom. Pri porovnávaní s mikroservisnou verziou tak slúži ako praktická ukážka toho, do akej miery je možné monolit „dotiahnuť“ smerom k dobre štruktúrovanému systému bez nutnosti jeho rozdelenia na samostatné služby.

\subsection{Požiadavky na riešenie}

Na základe analýzy domény hry Dáma a očakávaného spôsobu používania boli definované nasledujúce funkčné požiadavky:

\subsubsection{Funkčné požiadavky}

\begin{enumerate}
\item Základná hra Dáma – implementácia pravidiel hry (pohyb, povinné bije, kráľovanie, počítanie skóre).
\item Správa hráčov – zadanie niku a avatara pri začiatku hry, perzistentné ukladanie výsledkov.
\item Ukladanie výsledkov – persistencia výsledkov, komentárov a hodnotení v databáze.
\item Frontend rozhranie – React SPA aplikácia pre spustenie hry, zobrazenie dosky, zápis ťahov a leaderboard.
\item Lokálny režim – aplikácia musí podporovať lokálnu hru dvoch hráčov na jednom zariadení.
\end{enumerate}

\subsubsection{Nefunkčné požiadavky}

\begin{enumerate}
\item Výkon – rýchla odozva pri interakcii s hernou doskou, minimálne latencie pri vykresľovaní ťahov.
\item Udržiavateľnosť – kód má byť modulárny, dobre testovateľný a dokumentovaný.
\item Nasadzovanie a automatizácia – možnosť ľahkého spustenia lokálne cez Docker Compose a príprava na automatizované nasadzovanie.
\item Monitorovanie a observabilita predstavujú kľúčové nefunkčné požiadavky moderných distribuovaných systémov. Nástroje ako Prometheus a Grafana umožňujú sledovanie výkonnostných metrík, včasnú detekciu problémov a analýzu správania systému v produkčnom prostredí~\cite{PrometheusDocs2025,GrafanaDocs2025}.
\item Bezpečnosť – základné opatrenia na strane API (CORS, input validation) a princíp minimálnych oprávnení pri prístupe k databáze.
\item Rozšíriteľnosť a izolácia domén – architektúra musí umožniť jednoduché doplnenie ďalších komponentov a v prípade mikroservisov minimalizovať coupling medzi doménami.
\item Škálovateľnosť – možnosť horizontálneho rozširovania častí systému podľa potreby.
\item Spoľahlivosť – dostupnosť služby aj pri zlyhaní iných komponentov.
\end{enumerate}

Tieto nefunkčné požiadavky vychádzajú z cieľa porovnať architektúry nielen z hľadiska funkcionality, ale aj z prevádzkového správania a udržiavateľnosti~\cite{Newman2021,Richardson2018}.

\subsection{Návrh mikroservisnej architektúry}

Systém je rozdelený na samostatné služby: user-service – používateľské profily, nicky, avatary; game-service – celá herná logika vrátane validácie ťahov; score-service – záznam výsledkov; comment-service – komentáre; rating-service – hodnotenia; gateway-service – centrálny vstupný bod; frontend – React SPA.

Služby komunikujú synchrónne cez REST a využívajú spoločnú databázu PostgreSQL, pričom každá služba pracuje len s tabuľkami vlastnej domény.

Pri návrhu hraníc medzi jednotlivými službami sa vychádzalo z doménového členenia aplikácie a zo snahy minimalizovať závislosti medzi službami. Každá služba je zodpovedná za svoju časť domény a vystavuje rozhranie, ktoré je možné ďalej rozširovať bez nutnosti meniť vnútornú implementáciu iných služieb. Toto rozdelenie odráža odporúčania uvádzané v literatúre~\cite{Newman2021,Richardson2018}, podľa ktorých by mikroservisy mali rešpektovať hranice biznis domén a nemali by zdieľať vnútorné dátové štruktúry.

V kontexte hry Dáma to znamená napríklad oddelenie logiky priebehu hry od správy profilov hráčov a evidencie výsledkov. Takýto návrh uľahčuje nielen škálovanie konkrétnych častí systému, ale aj ich samostatný vývoj, testovanie a nasadzovanie.

\subsection{Komunikácia a konzistencia}

Monolit používa ACID transakcie v rámci jednej databázovej schémy; všetky operácie prebiehajú v rámci jedného procesu, čo zjednodušuje udržiavanie konzistencie.

V mikroservisnej architektúre je doménová logika oddelená do viacerých služieb. API kontrakty zabezpečujú stabilitu komunikácie a jasne definujú formáty dát.

Pri zápise výsledkov prebieha typická sekvencia:

frontend → gateway → game-service → score-service → databáza.

Výhodou mikroservisného prístupu je lepšia modularita a možnosť samostatne rozvíjať jednotlivé domény, nevýhodou je mierne zvýšená latencia a potreba riešiť konzistenciu v distribuovanom prostredí.

Komunikácia medzi službami
\begin{itemize}
\item Synchrónny REST – používa sa pre okamžité volania (napr. frontend → gateway → game-service).
\item Asynchrónne udalosti (teoreticky) – pomocou message brokeru (napr. RabbitMQ alebo Apache Kafka) pre decoupling a eventual consistency.
\item API kontrakty – každá služba definuje vlastné DTO, čím sa zabezpečuje nezávislosť a možnosť samostatného vývoja.
\end{itemize}

\subsubsection{Správa stavu hry}

Herné pole a logika hry sú spracovávané v pamäti aplikácie pomocou objektov. Po ukončení hry game-service zapíše výsledok do PostgreSQL (score, rating, comment).

\subsubsection{Konzistencia a transakcie}

Pri oddelení perzistentných služieb sa v mikroservisnej architektúre často uprednostňuje model eventual consistency, ktorý umožňuje vyššiu odolnosť systému a lepšiu škálovateľnosť za cenu dočasnej nekonzistencie dát~\cite{Richardson2018,KafkaDocs2025}.

Napríklad po ukončení hry game-service vyšle udalosť na score-service. Pre zložitejšie viackrokové transakcie sa v distribuovaných systémoch často používa SAGA pattern, ktorý umožňuje riadenie transakcií bez potreby globálnych ACID transakcií. Tento prístup je detailne popísaný v literatúre zameranej na návrh mikroservisných architektúr~\cite{Richardson2018}.

\subsubsection{Bezpečnosť a politika prístupu}

\begin{itemize}
\item gateway zodpovedá za základné bezpečnostné politiky (CORS, rate-limiting);
\item input validation je aplikovaná na každej vrstve API;
\item prístup do databáz je nastavený podľa princípu least privilege.
\end{itemize}

\subsection{Technologická analýza}

Na základe rámcového technologického kontextu uvedeného v predchádzajúcej podkapitole nasleduje detailnejšia analýza jednotlivých použitých nástrojov a ich vzťahu k nefunkčným požiadavkám systému.

\subsubsection{Hlavný technologický stack a jeho využitie}

\begin{itemize}
\item Spring Boot – framework pre backend: umožňuje rýchly vývoj REST API a spracovanie logiky mikroservisov. V projekte sa používa vo všetkých službách (user-service, game-service, score-service a ďalšie). Poskytuje modulárnu štruktúru, integráciu so Spring Data pre prácu s databázou, Spring Security pre autentifikáciu a autorizáciu a Spring Web pre tvorbu REST endpointov~\cite{SpringBootDocs2025,SpringWebRESTDocs2025}.
\item PostgreSQL – relačná databáza: slúži na perzistentné uchovávanie dát o hráčoch, skóre, komentároch a hodnoteniach. Používa sa pre doménové tabuľky jednotlivých služieb~\cite{PostgreSQLDocs2025}.
\item Docker a Docker Compose – kontajnerizácia: izolujú jednotlivé služby a uľahčujú lokálne spustenie a testovanie celého systému~\cite{Merkel2014Docker,DockerCompose2025}.
\item Spring Cloud Gateway – API Gateway: centralizované smerovanie požiadaviek, autentifikácia, CORS, rate-limiting~\cite{SpringCloudGateway2025}.
\item React – single-page aplikácia: React slúži ako frontend v oboch verziách systému. Zabezpečuje interaktívne používateľské rozhranie, renderovanie hernej dosky a komunikáciu cez REST API~\cite{ReactDocs2025}.
\end{itemize}

\subsection{Testovanie a meranie výkonnosti}

Testovanie zahŕňa jednotkové testy (JUnit 5 + Mockito), integračné testy (Spring Test + Testcontainers), validačné testy API kontraktov (Pact) a end-to-end scenáre simulujúce reálnu hru.

Výkonnostné testy zahŕňajú reakčnú dobu pri vykreslení ťahu, rýchlosť zápisu výsledku, spotrebu CPU/RAM počas ťahov a odozvu systému pri simultánnych hrách.

\subsubsection{Metodika merania výkonu}

Merania budú prebiehať v kontrolovanom prostredí (rovnaký hardvér, čisté inštancie). Scenáre zahŕňajú otvorenie hry, vykonanie všetkých možných ťahov, ukončenie hry a zápis výsledkov do databázy. Metriky zahŕňajú čas odozvy pri vykresľovaní ťahov a aktualizácii herného poľa, čas zápisu výsledkov a spotrebu CPU a pamäte počas hry. Použité nástroje zahŕňajú systémové metriky (Docker stats, htop).

Okrem kvantitatívnych ukazovateľov sa zhodnotí aj zložitosť údržby, čitateľnosť kódu a modularita.

\subsection{Zhrnutie analytickej časti a premostenie do návrhovej časti}

Analytická časť identifikovala hlavné obmedzenia monolitickej architektúry a navrhla moderné mikroservisné riešenie založené na Spring Boot, PostgreSQL, Docker, Spring Cloud Gateway a React.

Navrhnutá architektúra je modulárna, rozšíriteľná a prispôsobená pre objektívne porovnanie s monolitom podľa výkonových a údržbových metrík. Výsledky tejto analýzy budú tvoriť základ pre implementačnú časť práce, kde budú definované konkrétne API kontrakty, dátové modely, Docker Compose konfigurácie a testovacie scenáre.